We presented VL-KnG, a training-free framework that constructs persistent spatiotemporal knowledge graphs from monocular egocentric video, enabling structured and explainable scene understanding.
By decoupling knowledge construction from query processing, VL-KnG provides persistent memory and interpretable reasoning while avoiding the poor scaling of direct VLM inference with video duration.
Evaluation across three benchmarks shows that VL-KnG matches or surpasses frontier VLMs on embodied scene understanding tasks while achieving substantially lower query latency.
Graph-Enhanced Retrieval, combining structured subgraph reasoning with visual grounding, plays a key role in bridging the gap between text-level graph representations and frame-level localization.
Real-world robot deployment further demonstrates practical applicability, with constant-time query scaling independent of video length.
\textbf{Limitations.} Fine-grained spatial reasoning remains a limitation, as the knowledge graph encodes positions at the frame level without explicit 3D geometry. 
KG construction quality is also bounded by the underlying VLM's detection capability. 
Finally, the current STOA module assumes relatively static scenes where objects do not frequently change state or location, which may limit performance in highly dynamic environments.
Future work will explore dynamic environments with changing object states, integration of multi-modal sensing for richer graph construction, and scaling VL-KnG to long-duration video streams spanning hours of continuous observation. 
More broadly, our results suggest that persistent structured representations offer a promising path toward scalable and explainable embodied scene understanding beyond direct end-to-end VLM inference.